\subsection*{问“乘了几次”，而不是再发明一种运算口诀}
指数把次数送到倍率，对数反过来问某个正数对应多少次增长。$\log_2 8=3$ 的内容就是 $2^3=8$。当次数允许实数，对数成为连续的尺度转换，而不限于整数乘法次数。

若 $a=c^u,b=c^v$，乘积是 $c^{u+v}$，所以乘积的对数是两对数之和。这个法则直接来自指数的加法规则。商、幂的对数关系也可从同一来源推导，不必各自背一遍。

\subsection*{换底只是换了一把刻度尺}
底数二按翻倍计数，底数十按十倍计数，自然底数按连续增长计数。同一倍率在不同刻度中数值不同，彼此只差一个固定比例，因而出现换底公式。

输入为 $1/9$ 的三底对数是负二，负号表示反向增长或缩小；真数依旧为正。实对数的定义域不能因结果允许负数就扩展到负真数。底数一又无法区分任何次数，因此不能作为对数底。

对数尤其适合把很大的倍率范围压成可读的加性刻度，但必须知道图上的等距离表示等比例，而非等绝对差。后面的自然对数将把这条思想与微小相对变化联系起来。
